\section{משוואת הגלים בדו-ממד: הממברנה ומערכות כלליות}

\begin{mainbox}{מידע כללי}
\textbf{נושא:} משוואת הגלים בדו-ממד (ממברנה מלבנית), ניוון (Degeneracy), גאומטריה כללית והרמיטיות של הלפלסיאן.
\end{mainbox}

\subsection{הממברנה המלבנית: ניסוח ופתרון}

נדון בממברנה דו-ממדית (כגון דף או תוף מלבני) המתוחה במישור $xy$. הממברנה מקובעת בכל קצותיה.

\begin{definitionbox}{משוואת הגלים הדו-ממדית}
המשוואה המתארת את ההעתק $u(x,y,t)$ של הממברנה היא:
\begin{equation}
\frac{\partial^2 u}{\partial t^2} = c^2 \nabla^2 u = c^2 \left( \frac{\partial^2 u}{\partial x^2} + \frac{\partial^2 u}{\partial y^2} \right)
\end{equation}
כאשר $c = \sqrt{T/\sigma}$ היא מהירות הגל (כתלות במתיחות $T$ וצפיפות המסה המשטחית $\sigma$).
\end{definitionbox}

\subsubsection{תנאי שפה והתחלה}
הממברנה מוגדרת בתחום $0 \le x \le L_x$ ו-$0 \le y \le L_y$. הקצוות מקובעים (תנאי דיריכלה הומוגניים):
\begin{itemize}
    \item $u(0, y, t) = u(L_x, y, t) = 0$
    \item $u(x, 0, t) = u(x, L_y, t) = 0$
\end{itemize}
תנאי ההתחלה נתונים על ידי פונקציית הצורה ההתחלתית $f(x,y)$ והמהירות ההתחלתית $g(x,y)$.

\subsubsection{הפרדת משתנים}
נחפש פתרון בצורה של הפרדת משתנים מלאה:
\begin{equation}
u(x,y,t) = X(x)Y(y)T(t)
\end{equation}
הצבה במשוואה וחלוקה ב-$XYT$ מובילה ל:
\begin{equation}
\frac{1}{c^2} \frac{\ddot{T}}{T} = \frac{X''}{X} + \frac{Y''}{Y} = -k^2
\end{equation}
מכיוון שצד שמאל תלוי בזמן בלבד וצד ימין במרחב בלבד, שניהם שווים לקבוע משותף אותו נסמן $-k^2$. בנוסף, מכיוון ש-$X$ ו-$Y$ הם משתנים בלתי תלויים, כל אחד מהאיברים המרחביים חייב להיות קבוע בפני עצמו:
\begin{equation}
\frac{X''}{X} = -\lambda^2, \quad \frac{Y''}{Y} = -\mu^2
\end{equation}
כך שמתקיים הקשר (יחס הדיספרסיה):
\begin{equation}
k^2 = \lambda^2 + \mu^2
\end{equation}

\begin{resultbox}{הפונקציות העצמיות והערכים העצמיים}
פתרון המשוואות המרחביות עם תנאי השפה נותן:
\begin{equation}
X_n(x) = \sin\left( \frac{n\pi x}{L_x} \right), \quad n=1,2,3,\dots
\end{equation}
\begin{equation}
Y_m(y) = \sin\left( \frac{m\pi y}{L_y} \right), \quad m=1,2,3,\dots
\end{equation}
הפונקציה המרחבית המלאה (המוד, \LR{Mode}) היא:
\begin{equation}
\phi_{nm}(x,y) = \sin\left( \frac{n\pi x}{L_x} \right) \sin\left( \frac{m\pi y}{L_y} \right)
\end{equation}
\end{resultbox}

התדרים העצמיים של המערכת מתקבלים מהקשר $k_{nm}^2 = \lambda_n^2 + \mu_m^2$:
\begin{equation}
\omega_{nm} = c k_{nm} = c \sqrt{\left(\frac{n\pi}{L_x}\right)^2 + \left(\frac{m\pi}{L_y}\right)^2}
\end{equation}

החלק הזמני הוא פתרון של מתנד הרמוני:
\begin{equation}
T_{nm}(t) = A_{nm} \cos(\omega_{nm} t) + B_{nm} \sin(\omega_{nm} t)
\end{equation}

\subsection{ניוון (\LR{Degeneracy})}

במערכות רב-ממדיות, ייתכן מצב שבו לשתי פונקציות עצמיות שונות (צורות מרחביות שונות) יש את אותו ערך עצמי (אותו תדר). תופעה זו נקראת \textbf{ניוון}.

\begin{examplebox}{דוגמה: הממברנה הריבועית}
נניח כי הממברנה היא ריבוע, $L_x = L_y = L$. התדרים נתונים על ידי:
\begin{equation}
\omega_{nm} = \frac{c\pi}{L} \sqrt{n^2 + m^2}
\end{equation}
נתבונן במודים המעורבים:
\begin{itemize}
    \item עבור $(n,m) = (1,2)$: התדר הוא פרופורציונלי ל-$\sqrt{1^2+2^2} = \sqrt{5}$.
    \item עבור $(n,m) = (2,1)$: התדר הוא פרופורציונלי ל-$\sqrt{2^2+1^2} = \sqrt{5}$.
\end{itemize}
קיבלנו $\omega_{12} = \omega_{21}$, אך הפונקציות המרחביות שונות:
\[ \phi_{12} = \sin\left(\frac{\pi x}{L}\right)\sin\left(\frac{2\pi y}{L}\right) \neq \phi_{21} = \sin\left(\frac{2\pi x}{L}\right)\sin\left(\frac{\pi y}{L}\right) \]
\end{examplebox}

\begin{notebox}{קווי צומת (\LR{Nodal Lines})}
כאשר קיים ניוון, כל צירוף ליניארי של המודים המנוונים הוא גם פתרון עם אותו תדר.
\begin{equation}
u = A \phi_{12} + B \phi_{21}
\end{equation}
על ידי בחירת יחסים שונים בין $A$ ל-$B$, ניתן לייצר קווי צומת (האזורים בהם הממברנה לא זזה) בצורות גאומטריות מגוונות (אלכסונים, מעגלים משוערים וכו'), ולא רק קווים המקבילים לצירים.
\end{notebox}

\subsection{הפתרון הכללי ומציאת המקדמים}

מכיוון שמשוואת הגלים ליניארית, הפתרון הכללי הוא סופרפוזיציה של כל המודים האפשריים:
\begin{equation}
u(x,y,t) = \sum_{n=1}^{\infty} \sum_{m=1}^{\infty} \left[ A_{nm} \cos(\omega_{nm} t) + B_{nm} \sin(\omega_{nm} t) \right] \sin\left( \frac{n\pi x}{L_x} \right) \sin\left( \frac{m\pi y}{L_y} \right)
\end{equation}

כדי למצוא את המקדמים $A_{nm}$ ו-$B_{nm}$, נשתמש בתכונת האורתוגונליות של הפונקציות העצמיות בדו-ממד.
הפונקציות $\phi_{nm}(x,y)$ מהוות בסיס שלם ואורתוגונלי במלבן.

\begin{formulabox}{חישוב מקדמים (אינטגרל כפול)}
עבור תנאי ההתחלה $u(x,y,0) = f(x,y)$:
\begin{equation}
A_{nm} = \frac{4}{L_x L_y} \int_{0}^{L_x} \int_{0}^{L_y} f(x,y) \sin\left( \frac{n\pi x}{L_x} \right) \sin\left( \frac{m\pi y}{L_y} \right) \, dx \, dy
\end{equation}
עבור המהירות ההתחלתית $\dot{u}(x,y,0) = g(x,y)$:
\begin{equation}
B_{nm} = \frac{4}{L_x L_y \omega_{nm}} \int_{0}^{L_x} \int_{0}^{L_y} g(x,y) \sin\left( \frac{n\pi x}{L_x} \right) \sin\left( \frac{m\pi y}{L_y} \right) \, dx \, dy
\end{equation}
\end{formulabox}

\subsection{גאומטריה כללית והלפלסיאן}

כאשר צורת הממברנה אינה מלבנית, לא ניתן בהכרח להפריד משתנים בין הקואורדינטות המרחביות ($x$ ו-$y$). עם זאת, תמיד ניתן להפריד בין הזמן למרחב.

נניח פתרון מהצורה $u(\vec{r},t) = \psi(\vec{r})T(t)$. הצבה במשוואת הגלים מובילה ל:
\begin{itemize}
    \item משוואה זמנית (תמיד הרמונית): $\ddot{T} + \omega^2 T = 0$.
    \item משוואה מרחבית (\textbf{משוואת הלמהולץ}):
    \begin{equation}
    \nabla^2 \psi(\vec{r}) + k^2 \psi(\vec{r}) = 0
    \end{equation}
    כאשר $k = \omega/c$.
\end{itemize}
עלינו לפתור את משוואת הערכים העצמיים
\[
\nabla^2 \psi = -\lambda \psi \qquad (\lambda = k^2)
\]
בתחום $\Omega$ עם תנאי שפה על השפה $\Gamma$
(למשל תנאי דיריכלה $\psi|_{\Gamma} = 0$).

\subsubsection{הרמיטיות של הלפלסיאן}

כדי להבטיח שהפונקציות העצמיות $\psi_n$ פורשות את המרחב (סט שלם)
ושהערכים העצמיים $\lambda_n$ ממשיים, עלינו להוכיח שאופרטור הלפלסיאן
הוא אופרטור הרמיטי (צמוד לעצמו) תחת מכפלה פנימית המוגדרת על־ידי
אינטגרל:

\[
\langle f, g \rangle
=
\iint_{\Omega} f(\vec{r})\, g(\vec{r}) \, dS .
\]

\begin{theorembox}{משפט: הרמיטיות של $\nabla^2$ בתחום חסום}
בהינתן שתי פונקציות $u, v$ המקיימות תנאי שפה הומוגניים
(דיריכלה $u=0$ או נוימן $\frac{\partial u}{\partial n}=0$ על השפה),
מתקיים:
\begin{equation}
\iint_{\Omega} u \, (\nabla^2 v) \, dS
=
\iint_{\Omega} v \, (\nabla^2 u) \, dS .
\end{equation}
\end{theorembox}


נשתמש בזהות וקטורית עבור הדיברגנץ של שדה מהצורה $u \nabla v$:
\[
\nabla \cdot (u \nabla v)
=
\nabla u \cdot \nabla v
+
u \nabla^2 v .
\]
נבצע אינטגרציה על התחום $\Omega$ ונשתמש ב\textbf{משפט הדיברגנץ (גאוס)}
בדו־ממד:
\[
\iint_{\Omega}
\left(
u \nabla^2 v + \nabla u \cdot \nabla v
\right)
\, dS
=
\oint_{\Gamma}
(u \nabla v) \cdot \hat{n} \, dl
=
\oint_{\Gamma}
u \frac{\partial v}{\partial n} \, dl .
\]
נחליף את התפקידים של $u$ ו־$v$ ונחסר את שתי המשוואות
(הזהות השנייה של גרין):
\begin{equation}
\iint_{\Omega}
\left(
u \nabla^2 v - v \nabla^2 u
\right)
\, dS
=
\oint_{\Gamma}
\left(
u \frac{\partial v}{\partial n}
-
v \frac{\partial u}{\partial n}
\right)
\, dl .
\end{equation}
מאחר ש־$u$ ו־$v$ מקיימות תנאי שפה הומוגניים,
האינטגרל השפתי (אגף ימין) מתאפס.
לדוגמה, בתנאי דיריכלה:
\[
u|_{\Gamma} = v|_{\Gamma} = 0 ,
\]
ולכן האגף הימני שווה לאפס.
מכאן נובע שוויון בין האינטגרלים המרחביים,
ולפיכך אופרטור הלפלסיאן הוא אופרטור הרמיטי.
\begin{summarybox}{מסקנות למערכות כלליות}
\begin{enumerate}
    \item הערכים העצמיים $k_n^2$ הם ממשיים
    (ובדרך כלל חיוביים עבור משוואת הגלים).
    \item הפונקציות העצמיות
    $\psi_n(\vec{r})$
    המתאימות לערכים עצמיים שונים הן
    \textbf{אורתוגונליות}.
    \item אוסף הפונקציות העצמיות מהווה
    \textbf{בסיס שלם},
    המאפשר לפרוס כל פונקציה (ובפרט את תנאי ההתחלה)
    כטור:
    \[
    f(\vec{r}) = \sum_n A_n \psi_n(\vec{r}) .
    \]
\end{enumerate}
\end{summarybox}
